\section{Merging equal-colored triangles in a cell}
\label{sec:area_law_fan_runs}

Fix a fan in one dyadic square $Q$ and assign each triangle one of the
two colors $0,1$. The merging below applies to every such assignment;
see \cite[Section~11]{OpenAI2026AreaLaw}.

% Source: 10-geometry.tex, lines 313--318. This is the within-cell
% merging step. The assignment across actual opposing regions is separate.

\begin{definition}[Runs and their regions]
    \label{def:area_law_fan_runs}
    \lean{TNLean.PEPS.AreaLaw.Geometry.cellFanRunGraph,
        TNLean.PEPS.AreaLaw.Geometry.CellFanRun,
        TNLean.PEPS.AreaLaw.Geometry.cellFanRunRegion}
    \leanok
    \uses{def:area_law_cell_fan}
    Write $[a_e,b_e]$ for the elementary perimeter segment of triangle
    $P_e$, oriented counterclockwise, and let $\lambda(e)\in\{0,1\}$
    be its color. Two distinct triangles are adjacent when
    $b_e=a_{e'}$ or $b_{e'}=a_e$. Retain an adjacency precisely when
    the two colors agree. A run is a connected component of this
    retained graph. Its region is
    \begin{align}
        T_R=\bigcup_{e\in R}P_e.
    \end{align}
    Distinct components remain distinct identifiers, including when
    their colors agree.
\end{definition}

\begin{theorem}[Color constancy on a run]
    \label{thm:area_law_fan_run_color}
    \lean{TNLean.PEPS.AreaLaw.Geometry.cellFanRun_color_eq}
    \leanok
    \uses{def:area_law_fan_runs}
    Any two triangles in the same run have the same color.
\end{theorem}
\begin{proof}\leanok
    \uses{def:area_law_fan_runs}
    Along every retained edge the colors agree. Following a finite
    path in a component gives the assertion.
\end{proof}

\begin{theorem}[Exact cover by run regions]
    \label{thm:area_law_fan_run_cover}
    \lean{TNLean.PEPS.AreaLaw.Geometry.cellFanRunRegions_cover}
    \leanok
    \uses{def:area_law_fan_runs}
    The run regions cover exactly the closed square:
    \begin{align}
        Q=\bigcup_R T_R.
    \end{align}
\end{theorem}
\begin{proof}\leanok
    \uses{def:area_law_fan_runs,thm:area_law_cell_fan_cover}
    Every elementary segment belongs to one component. Thus the
    union of all run regions is the union of all fan triangles, which
    is $Q$.
\end{proof}

\begin{theorem}[Colors across adjacent distinct runs]
    \label{thm:area_law_fan_run_opposite}
    \lean{TNLean.PEPS.AreaLaw.Geometry.cellFanRun_adjacent_colors_ne}
    \leanok
    \uses{def:area_law_fan_runs}
    Suppose that two adjacent fan triangles belong to distinct runs.
    Their colors are different, and hence are opposite members of
    $\{0,1\}$.
\end{theorem}
\begin{proof}\leanok
    \uses{def:area_law_fan_runs}
    If their colors agreed, their adjacency would be retained and
    they would lie in the same component.
\end{proof}

\begin{theorem}[A monochromatic fan gives one whole-cell region]
    \label{thm:area_law_fan_run_monochromatic}
    \lean{TNLean.PEPS.AreaLaw.Geometry.cellFanRun_all_equal}
    \leanok
    \uses{def:area_law_fan_runs}
    If every triangle has the same color, all triangles have the same
    run identifier, and every run region is the whole closed square.
\end{theorem}
\begin{proof}\leanok
    \uses{def:area_law_fan_runs,thm:area_law_cell_fan_cover}
    The one or two elementary segments on each side are connected
    in their perimeter order. The last segment of each side is
    adjacent to the first segment of the next side, including the
    fourth-to-first transition. Under constant coloring all these
    adjacencies are retained, so the graph is connected. Its single
    component contains every triangle, whose union is $Q$.
\end{proof}
